\documentclass[11pt]{article}

\usepackage[margin=1in]{geometry}
\usepackage[T1]{fontenc}
\usepackage[utf8]{inputenc}
\usepackage{lmodern}
\usepackage{microtype}
\usepackage{amsmath,amssymb,amsthm,mathtools}
\usepackage{booktabs}
\usepackage{tabularx}
\usepackage{array}
\usepackage{enumitem}
\usepackage{xcolor}
\usepackage[colorlinks=true,linkcolor=blue!50!black,citecolor=blue!50!black,urlcolor=blue!50!black]{hyperref}

\setlist{itemsep=2pt,topsep=4pt}
\newcommand{\Ltwo}{L\textsuperscript{2}}
\newcolumntype{L}{>{\raggedright\arraybackslash}X}
\newcommand{\KL}{D_{\mathrm{KL}}}
\newcommand{\E}{\mathbb{E}}
\newcommand{\R}{\mathbb{R}}
\newcommand{\vu}{{v\mid u}}
\newcommand{\uv}{{u\mid v}}
\newcommand{\cS}{\mathcal{S}}

\theoremstyle{plain}
\newtheorem{proposition}{Proposition}
\newtheorem{corollary}[proposition]{Corollary}
\newtheorem{lemma}[proposition]{Lemma}
\theoremstyle{definition}
\newtheorem{definition}[proposition]{Definition}
\theoremstyle{remark}
\newtheorem{remark}[proposition]{Remark}

\title{Bounded scores and well-specified families:\\
when the paper's contrastive objective works, and when the loss variant matters}
\author{}
\date{}

\begin{document}
\maketitle

\begin{abstract}
\noindent
The goals~1--3 write-up (\texttt{docs/goals1-3\_writeup.pdf}) ended with two empirical
patterns. First, every catastrophic failure was an unbounded score on heavy-tailed or spiky
data. Second, the paper's sharpest predictions (the dot product's wrong spread, P1, and the
joint loss's shrinkage, P4) held for linear encoders and disappeared for MLPs. This note gives
both patterns a theoretical footing inside the paper's own framework~\cite{bst}.
\emph{Part~I (boundedness):} a score that is bounded above defines a valid model for every
data distribution, and its bank forecast is consistent. A score that is linear in the future
does not: on Student-$t$ futures its model does not exist for any informative parameter, ReLU
MLPs do not change this, and the bank forecast provably collapses onto a single future as
the bank grows. \emph{Part~II (misspecification):} the one-way, two-way and joint losses
have the same population target, the log density ratio $\log r$, up to different free terms.
So they agree whenever the score family can represent that target, and can differ only when
it cannot. In the Gaussian case this recovers P1, P3 and P4 as instances of misspecification.
It also gives a new result: linear \Ltwo{} trained two-way or joint is misspecified on any
correlated Gaussian. At weak dependence it loses up to $3/8$ of the available information,
yet only $0.3\%$ on the goal-1 field, which is why goal~1 did not see it. Every numerical
claim is reproduced by \texttt{test\_bed/theory/misspecification.py}.
\end{abstract}

%======================================================================
\section{Setting and notation}

We follow the paper~\cite{bst}. The past is $u\in U$, the future is $v\in V$, the data
distribution is $\mu(du,dv)$, and its marginals are $\mu_u,\mu_v$. Throughout we assume $\mu$ and
$\mu_u\otimes\mu_v$ are mutually absolutely continuous, and write
\[
  r(u,v) = \frac{d\mu}{d(\mu_u\otimes\mu_v)}(u,v), \qquad
  \frac{d\mu_\vu(\cdot\mid u)}{d\mu_v}(v) = r(u,v).
\]
So $\log r$ is the pointwise mutual information, and
$I(u;v) = \E_\mu[\log r] = \E_u\,\KL(\mu_\vu\,\|\,\mu_v)$.

A trained model is a \emph{score} $s:U\times V\to\R$. It defines a joint model and its
conditionals (the paper's eqs.~(13)--(16)):
\[
  \nu(du,dv) = \frac{e^{s(u,v)}}{Z}\,\mu_u(du)\mu_v(dv),\qquad
  \nu_\vu(dv\mid u) = \frac{e^{s(u,v)}}{Z(u)}\,\mu_v(dv),
\]
\[
  Z = \iint e^{s}\,d\mu_u\,d\mu_v,\qquad Z(u) = \int e^{s(u,v)}\,\mu_v(dv).
\]
Up to constants that do not depend on $s$, the losses we trained are (the paper's
Thm~3.2, Thm~3.8 and Remark~3.12):
\begin{align*}
  J_\vu(s)   &= \E_{u}\,\KL\big(\mu_\vu(\cdot\mid u)\,\|\,\nu_\vu(\cdot\mid u)\big)
               && \text{one-way (forecasting direction)},\\
  J_{2}(s)   &= \lambda_u J_\vu(s) + \lambda_v J_\uv(s),\ \ \lambda_u,\lambda_v>0
               && \text{two-way},\\
  J_{\rm joint}(s) &= \KL(\mu\,\|\,\nu)
               && \text{joint}.
\end{align*}
The forecast we score is the paper's eq.~(31) run on a finite \emph{bank} $v_1,\dots,v_N$ of
real futures. This is self-normalised importance sampling with proposal $\mu_v$:
\[
  \nu^N_\vu(\cdot\mid u) = \sum_{j=1}^N w_j\,\delta_{v_j},\qquad
  w_j = \frac{e^{s(u,v_j)}}{\sum_k e^{s(u,v_k)}}, \qquad
  \mathrm{ESS} = \Big(\sum_j w_j^2\Big)^{-1}.
\]

\paragraph{Score families.} With encoders $f$ (past) and $g$ (future) and a scale $\kappa>0$
(the logit scale), the three families in \texttt{test\_bed/losses/clip\_losses.py} are
\[
  s_{\rm dot} = \kappa\,\langle f(u),g(v)\rangle,\qquad
  s_{\rm ncos} = \kappa\,\Big\langle \tfrac{f(u)}{|f(u)|},\tfrac{g(v)}{|g(v)|}\Big\rangle,\qquad
  s_{\Ltwo} = -\tfrac{\kappa}{2}\,|f(u)-g(v)|^2 .
\]
So $s_{\Ltwo}\le 0$ and $|s_{\rm ncos}|\le\kappa$, while $s_{\rm dot}$ is unbounded whenever
$g$ is.

%======================================================================
\section{Part I: boundedness decides whether the model exists}

\begin{definition}
A score $s$ is \emph{admissible} for $\mu$ (in the forecasting direction) if
$Z(u)<\infty$ for $\mu_u$-a.e.\ $u$. It is \emph{jointly admissible} if $Z<\infty$.
\end{definition}

If $s$ is not admissible, then $\nu_\vu(\cdot\mid u)$ does not exist and $J_\vu(s)=+\infty$.
The probabilistic reading of contrastive learning then has nothing to converge to.

\begin{proposition}[Bounded above $\Rightarrow$ always admissible]\label{prop:bounded}
If $s\le c$, then $Z(u)\le e^c$ and $Z\le e^c$ for \emph{every} $\mu$. In particular \Ltwo{}
($c=0$) and normalised cosine ($c=\kappa$) are admissible for every data distribution, with
any encoders.
\end{proposition}

\begin{proposition}[Scores linear in the future need exponential moments]\label{prop:linear}
Let $s(u,v)=\langle a(u),Hv\rangle+b(u)$, which covers the unnormalised dot product with a
linear future encoder. Then $Z(u)=e^{b(u)}M_{\mu_v}(H^\top a(u))$, where $M_{\mu_v}$ is the moment
generating function of $\mu_v$. If $\mu_v$ is a multivariate Student-$t$ (any degrees of
freedom), then $M_{\mu_v}(w)=\infty$ for every $w\neq 0$. So $s$ is admissible only if
$H^\top a(u)=0$, that is, only if the model ignores the past. Every informative parameter has
$J_\vu=+\infty$.
\end{proposition}

\begin{proposition}[ReLU MLPs do not rescue the dot product]\label{prop:relu}
Let $g$ be a ReLU network with biases and let $g_\infty$ be the same network with its biases
removed. Then $g_\infty$ is continuous and positively homogeneous, and $\sup_v|g(v)-g_\infty(v)|<\infty$.
If $\langle f(u),g_\infty(v_0)\rangle>0$ for some $v_0$ and $\mu_v$ is an elliptical Student-$t$,
then $Z(u)=\infty$. So the dot product is admissible at $u$ only if
$\langle f(u),g_\infty(v)\rangle\le 0$ for all $v$: the score may not grow in any direction of
the future. Nothing in training enforces this. A bounded output layer for $g$ (e.g.\ $\tanh$)
does enforce it, by Proposition~\ref{prop:bounded}, without normalising anything.
\end{proposition}

Admissibility is not just a population-level nicety. It is exactly the condition under which
the bank forecast works.

\begin{proposition}[Admissible $\Leftrightarrow$ the bank forecast is consistent]\label{prop:snis}
Fix $u$ and draw the bank i.i.d.\ from $\mu_v$. If $Z(u)<\infty$, then
$\nu^N_\vu(\cdot\mid u)\to\nu_\vu(\cdot\mid u)$ weakly, almost surely, and
\[
  \frac{\mathrm{ESS}}{N}\;\to\;\frac{1}{\E_{\mu_v}\big[(d\nu_\vu/d\mu_v)^2\big]}
  \quad\text{whenever that second moment is finite.}
\]
If $Z(u)=\infty$, the normaliser $\frac1N\sum_j e^{s(u,v_j)}$ diverges and there is no limit
to converge to.
\end{proposition}

When $Z(u)=\infty$ we can say what happens instead. The forecast collapses onto one future.

\begin{proposition}[Collapse under heavy-tailed scores]\label{prop:collapse}
Fix $u$ and let $X_j=s(u,v_j)$ be i.i.d.\ with a regularly varying right tail,
$P(X>x)=x^{-\alpha}L(x)$ with $\alpha>0$ and $L$ slowly varying. Then
\[
  \max_j w_j \;\to\; 1 \quad\text{in probability as } N\to\infty .
\]
For the linear score of Proposition~\ref{prop:linear} on Student-$t$ futures with $\nu$ degrees
of freedom, $X$ is a scaled univariate $t$, so this holds with $\alpha=\nu$.
\end{proposition}

So for an unbounded score on heavy tails, \emph{a bigger bank makes the forecast worse}. It
supplies more extreme futures to collapse onto. For bounded scores a bigger bank only helps
(Proposition~\ref{prop:snis}). Table~\ref{tab:collapse} checks both directions: a linear
score with $a=1$ on unit-variance futures, against an \Ltwo{} score on the same Student-$t$
futures. For light tails the limit is explicit: with Gaussian futures and a linear score,
$\mathrm{ESS}/N\to e^{-\sigma^2}$ with $\sigma^2=\mathrm{Var}\,s(u,v)$. Here that is
$e^{-1}=0.368$, as observed.

\begin{table}[h]
\centering\small
\begin{tabular}{@{}rccc@{}}
\toprule
bank size $N$ & linear score, Gaussian & linear score, Student-$t$ ($\nu=3$) & \Ltwo{} score, Student-$t$ ($\nu=3$)\\
\midrule
$10^2$ & 0.081 \;(0.41) & 0.240 \;(0.30)\phantom{00} & 0.018 \;(0.79)\\
$10^3$ & 0.016 \;(0.38) & 0.455 \;(0.038)\phantom{0} & 0.002 \;(0.79)\\
$10^4$ & 0.003 \;(0.37) & 0.760 \;(0.0007) & $<10^{-3}$ (0.79)\\
$10^5$ & 0.001 \;(0.37) & 0.902 \;($10^{-5}$)\phantom{0} & $<10^{-4}$ (0.79)\\
$10^6$ & $<10^{-4}$ (0.37) & 0.954 \;($10^{-6}$)\phantom{0} & $<10^{-5}$ (0.79)\\
\bottomrule
\end{tabular}
\caption{Largest bank weight, mean over 200 banks, with $\mathrm{ESS}/N$ in brackets. Only
the unbounded score on heavy tails collapses, and it collapses more as the bank grows.}
\label{tab:collapse}
\end{table}

\begin{proposition}[Bounded on both sides: guaranteed efficiency, capped sharpness]\label{prop:twosided}
If $|s|\le\kappa$ (normalised cosine), then for every $u$
\[
  e^{-2\kappa}\le\frac{d\nu_\vu}{d\mu_v}\le e^{2\kappa},\qquad
  \KL(\nu_\vu\,\|\,\mu_v)\le 2\kappa,\qquad
  \frac{\mathrm{ESS}}{N}\to\text{a limit}\ \ge e^{-2\kappa}.
\]
So normalised cosine has a bank requirement that is uniform in $u$. The price is that no
forecast can be more than $e^{2\kappa}$ times as concentrated as climatology, and none can
carry more than $2\kappa$ nats of information about the future. The sharpness of the whole
model is set by a single number, the logit scale.
\end{proposition}

\begin{table}[h]
\centering\small
\begin{tabularx}{\textwidth}{@{}lLLL@{}}
\toprule
& unbounded (dot product) & bounded above (\Ltwo{}) & bounded both sides (norm.\ cosine) \\
\midrule
model exists & only with light-tailed futures & always & always \\
bank forecast consistent & only with light-tailed futures & always & always \\
bank needed & can be infinite (collapse) & finite, $\approx e^{\KL(\nu_\vu\|\mu_v)}$~\cite{cd}, can be large & $\le$ order $e^{2\kappa}$, uniform in $u$ \\
sharpness & unlimited & unlimited & capped by $\kappa$ \\
\bottomrule
\end{tabularx}
\caption{The boundedness hierarchy.}
\label{tab:hierarchy}
\end{table}

\paragraph{What this adds to the write-up's \S7.} The write-up guessed that $Z(u)$ is a moment
generating function and that bounded scores always define a valid model.
Propositions~\ref{prop:bounded}--\ref{prop:linear} make that precise. Proposition~\ref{prop:relu}
explains why MLP cosine still failed on heavy tails in goal~2, and
Proposition~\ref{prop:snis} turns validity of the population model into consistency of the
forecast we actually score. Proposition~\ref{prop:collapse} turns ``a few extreme futures take
all the weight'' into a theorem. One honesty point: real rain is physically bounded, so
there $Z(u)$ is finite and Proposition~\ref{prop:collapse} does not apply literally. What
applies is its finite-$N$ shadow: $\mathrm{ESS}/N\approx 1/\E[(d\nu/d\mu_v)^2]$ is tiny when the
score's distribution over the bank is heavy-tailed. The direct check proposed in the
write-up (embedding norm against rain amount, and weight concentration) is still the way
to confirm the mechanism on rain.

%======================================================================
\section{Part II: misspecification decides when the loss variant matters}

All three losses aim at the same object, the log density ratio, but they leave different
parts of it free.

\begin{proposition}[Common target]\label{prop:target}
For any score $s$:
\begin{enumerate}[label=(\roman*)]
\item $J_\vu(s)=0 \iff s(u,v)=\log r(u,v)+c(u)$ for some function $c$;
\item $J_\uv(s)=0 \iff s(u,v)=\log r(u,v)+c'(v)$ for some function $c'$;
\item $J_2(s)=0 \iff J_{\rm joint}(s)=0 \iff s=\log r+\text{const}$.
\end{enumerate}
\end{proposition}

\begin{definition}
A score family $\cS$ is \emph{well-specified} for a loss if that loss attains $0$ on $\cS$.
By Proposition~\ref{prop:target}, it is well-specified for the one-way loss if it contains
some $\log r+c(u)$, and for the two-way and joint losses if it contains $\log r+\text{const}$.
\end{definition}

\begin{corollary}[Agreement under well-specification]\label{cor:agree}
If $\cS$ contains $\log r+\text{const}$, then the population minimisers of the one-way, the
two-way (for any $\lambda_u,\lambda_v>0$) and the joint loss all give exactly
$\nu_\vu=\mu_\vu$. The choice of loss variant cannot matter at the population level.
\end{corollary}

So a difference between loss variants is a symptom of misspecification. The converse is
false: P3 (one-way and two-way agree for the dot product) is a case of a misspecified family
whose variants agree anyway. The paper's Theorem~5.1 shows that for that family the
weighting only rescales the loss.

\subsection{The Gaussian case: P1, P3 and P4 as misspecification}

Let $\mu$ be a centred Gaussian with covariance blocks $C_{uu},C_{uv},C_{vv}$, joint
precision $P$, and $D=\mathrm{blockdiag}(C_{uu},C_{vv})$. Then
\[
  \log r(u,v) = -\tfrac12 u^\top(P_{uu}-C_{uu}^{-1})u \;-\; u^\top P_{uv} v
                \;-\; \tfrac12 v^\top(P_{vv}-C_{vv}^{-1})v + \text{const},
\]
with $P_{vv}=C_{v|u}^{-1}$ and $P_{uu}=C_{u|v}^{-1}$, so both ``self'' terms are positive
semidefinite. With linear encoders every score is a quadratic,
\[
  s(u,v) = -\tfrac12 u^\top R u + u^\top A v - \tfrac12 v^\top S v,
\]
and the families differ only in which $(R,A,S)$ they allow:
\begin{center}\small
\begin{tabular}{@{}lll@{}}
\toprule
family & allowed $(R,A,S)$ & well-specified for \\
\midrule
dot product & $R=S=0$ & nothing (if $C_{uv}\neq0$) \\
\Ltwo{} & $\begin{psmallmatrix}R&-A\\-A^\top&S\end{psmallmatrix}=[G,-H]^\top[G,-H]\succeq0$ & one-way only \\
dot $+$ quadratics & $R,S$ symmetric, $A$ free & every loss \\
\bottomrule
\end{tabular}
\end{center}

\begin{itemize}
\item \textbf{Dot product (P1).} It has no $v$-only term, so $\nu_\vu$ has covariance
  $C_{vv}$ whatever $A$ is. The mean can be right, but the spread is stuck at climatology,
  which is the paper's Corollary~5.2.
\item \textbf{Dot product, joint (P4).} It is additionally missing the $u$-only term. And
  the joint normaliser is finite only if $\|C_{uu}^{1/2}AC_{vv}^{1/2}\|<1$, which is where the
  paper's shrinkage function $h$ (Thm~5.6) comes from. Our numerical minimiser reproduces $h$
  exactly.
\item \textbf{\Ltwo{}, one-way (P2).} It is well-specified: the $v$-only term
  $-\frac12 v^\top S v$ with $S\succeq 0$ can match $P_{vv}-C_{vv}^{-1}$. This is the paper's
  Theorem~5.3 and Corollary~5.4.
\end{itemize}

\begin{proposition}[Linear \Ltwo{} is misspecified for the two-way and joint losses]\label{prop:l2joint}
Under linear \Ltwo{}, the model's joint precision is $D^{-1}+Q$ with $Q\succeq0$, so every
model joint covariance satisfies $(D^{-1}+Q)^{-1}\preceq D$. The data covariance satisfies
$C\not\preceq D$ whenever $C_{uv}\ne0$, because $D-C$ has eigenvalues $\pm$ the singular values
of $C_{uv}$. Hence $\log r+\text{const}$ is not in the family, and neither the two-way nor the
joint loss can be minimised to zero. In words: an \Ltwo{} model can couple $u$ and $v$ only by
shrinking both marginals, and real data couple them without shrinking.
\end{proposition}

The PSD constraint also ties the two directions together. By a Schur complement,
$S\succeq A^\top R^{+}A$: a strong past-to-future coupling $A$ forces either a sharp forecast
(large $S$) or a sharp backward conditional (large $R$). We return to this in \S\ref{sec:tension}.

\begin{proposition}[Scalar closed form]\label{prop:scalar}
For scalar $u,v$ with unit variances and correlation $\rho\in(0,1)$, the joint-\Ltwo{} optimum
forecasts
\[
  \nu_\vu(\cdot\mid u) = \mathcal N\Big(\frac{\rho}{2-\rho}\,u,\ \frac{2(1-\rho)}{2-\rho}\Big),
  \qquad\text{against the truth}\qquad
  \mathcal N\big(\rho\,u,\ 1-\rho^2\big).
\]
The mean is shrunk by a factor $1/(2-\rho)$ and the forecast is over-confident (its variance
is too small by a factor $2/(2+\rho-\rho^2)$). The fraction of the available information
lost, $J_\vu/I(u;v)$, tends to $3/8$ as $\rho\to0$ and to $0$ as $\rho\to1$.
\end{proposition}

Table~\ref{tab:miss} gives the full grid at the population optimum of each loss, as the
fraction of $I(u;v)$ that the resulting $v\mid u$ forecast loses. It is computed in closed
form from Gaussian KLs, with no training and no sampling.

\begin{table}[h]
\centering\small
\begin{tabular}{@{}lc|ccc|ccc|c@{}}
\toprule
& & \multicolumn{3}{c|}{dot product} & \multicolumn{3}{c|}{\Ltwo{}} & dot $+$ quad. \\
setting & $I(u;v)$ & 1-way & 2-way & joint & 1-way & 2-way & joint & all three \\
\midrule
$\rho=0.1$ & 0.005 & 0.005 & 0.005 & 0.005 & 0 & 0.294 & 0.333 & 0 \\
$\rho=0.3$ & 0.047 & 0.046 & 0.046 & 0.051 & 0 & 0.219 & 0.251 & 0 \\
$\rho=0.5$ & 0.144 & 0.131 & 0.131 & 0.157 & 0 & 0.148 & 0.170 & 0 \\
$\rho=0.7$ & 0.337 & 0.272 & 0.272 & 0.323 & 0 & 0.081 & 0.092 & 0 \\
$\rho=0.9$ & 0.830 & 0.512 & 0.512 & 0.571 & 0 & 0.021 & 0.023 & 0 \\
goal-1 GP, $50\to50$ & 1.627 & 0.704 & 0.704 & 0.746 & 0 & 0.003 & 0.003 & 0 \\
\bottomrule
\end{tabular}
\caption{Fraction of the information $I(u;v)$ (nats) lost by the forecast at the population
optimum: 0 is perfect, 1 is climatology. The last row uses the exact covariances of the
goal-1 Gaussian domain (\texttt{test\_bed/config.yaml}). Source:
\texttt{test\_bed/results/theory/misspecification.md}.}
\label{tab:miss}
\end{table}

Three readings of Table~\ref{tab:miss}:
\begin{itemize}
\item The dot product and two-way/joint \Ltwo{} fail in \emph{opposite regimes}. The dot
  product's error is in the spread, which is small when dependence is weak and large when it
  is strong. \Ltwo{}'s two-way/joint error is the coupling-by-shrinking constraint, which
  bites hardest when dependence is weak. At $\rho=0.1$, two-way \Ltwo{} loses 29\% of the
  information and the dot product loses 0.5\%.
\item The goal-1 field is in the strong-dependence regime ($I=1.63$ nats). There, two-way and
  joint \Ltwo{} lose 0.3\%. That is why goal~1 measured one-way, two-way and joint \Ltwo{}
  all at $+0.0\%$ on the Gaussian field and why the write-up's scorecard gives P3 a tick for \Ltwo{}. The theory
  says that tick is regime-dependent.
\item Adding free $u$-only and $v$-only quadratic terms to the dot product makes every loss
  exact. This sharpens the write-up's open question~1: a future-only term should fix P1 for
  the one-way loss, and adding a past-only term as well should make P3 and P4 disappear.
\end{itemize}

\subsection{Nonlinear encoders: well-specified in the bulk, not in the tails}

If $f$ and $g$ can output a constant coordinate (biases allow this), then
$\langle f(u),g(v)\rangle$ contains terms $c(u)+d(v)$. Sums $\sum_k f_k(u)g_k(v)$ are dense in
the continuous functions on compact sets (Stone--Weierstrass). So with enough width, the MLP
dot-product family is approximately well-specified \emph{on compacts}. By
Corollary~\ref{cor:agree}, the loss variants then converge, and P1 and P4 fade. That is
goal~2.

Approximation on compacts says nothing about growth at infinity, which is exactly what
Part~I is about. The two parts therefore divide the work cleanly:
\begin{center}
\emph{in the bulk of the data, misspecification decides; in the tails, boundedness decides.}
\end{center}
Linear normalised cosine is the remaining misspecified case. It sees the past only through
the direction $Au/|Au|$, so it cannot represent a $\log r$ that depends on the level of the
past. That explains its $+12$--$23\%$ with linear encoders (write-up, Table~2).

%======================================================================
\section{Where the theory and the experiments disagree}\label{sec:tension}

On weather, the population theory predicts the wrong ranking among the \Ltwo{} variants. One-way
\Ltwo{} is the well-specified one, yet it lost to joint in every case, and to two-way in two of three:
\begin{center}\small
\begin{tabular}{@{}lccc@{}}
\toprule
skill (2 seeds) & \Ltwo{} 1-way & \Ltwo{} 2-way & \Ltwo{} joint \\
\midrule
rain, linear (vs climatology) & $-0.106$ & $+0.011$ & $+0.012$ \\
temperature, MLP (vs persistence) & $+0.152$ & $+0.243$ & $+0.280$ \\
temperature, linear (vs persistence) & $+0.352$ & $+0.344$ & $+0.383$ \\
\bottomrule
\end{tabular}
\end{center}
The population argument bounds what joint and two-way can lose: at most about $3/8$ (joint)
or $1/3$ (two-way) of $I(u;v)$ in the scalar Gaussian case. On rain, $I(u;v)$ itself is tiny
(the best model beats climatology by 1\%), so that bound is negligible in absolute terms. The
observed gaps are far larger, so they must be finite-sample effects. A consistent reading:
\begin{itemize}
\item One-way \Ltwo{} has no check on sharpness. On rain its bank weights collapsed (minimum
  ESS 1.6 linear, median ESS 3 with MLPs; \texttt{summary\_B.md}), even though the score is
  bounded above. A bounded-above score still allows an arbitrarily sharp forecast
  (Table~\ref{tab:hierarchy}).
\item Two-way and joint \Ltwo{} do have such a check. Through $S\succeq A^\top R^{+}A$, a sharp
  forecast must come with a sharp backward conditional, which the two-way loss scores. The
  joint loss also scores the model's marginals, which a collapsed forecast distorts.
\end{itemize}
This is a bias--variance reading of the write-up's open question~3 (``why is joint
consistently best?''). The misspecification is a bias that is small exactly when the signal
is weak, and it buys a regulariser against over-sharpening, which is worst exactly then. It
predicts that the one-way deficit shrinks as training data grows, and turns into the
population ranking in the large-data limit. Goals 1--2, with plentiful synthetic data, are
consistent with that: there all \Ltwo{} variants tie. This reading is a hypothesis, not a
result.

%======================================================================
\section{Predictions and cheap tests}

Each item below is a direct consequence of a proposition above and can be run with the
existing test bed.
\begin{enumerate}
\item \textbf{Bank-size sweep} (Prop.~\ref{prop:collapse}). For unnormalised cosine on the
  df${}=3$ domain, the largest weight should rise and skill should fall as the bank grows. For
  \Ltwo{} and normalised cosine, the opposite.
\item \textbf{Bounded, unnormalised future encoder} (Props.~\ref{prop:bounded},
  \ref{prop:relu}). The dot product with a $\tanh$ output on $g$ should not collapse on heavy
  tails. This separates ``bounded'' from ``normalised''.
\item \textbf{Quadratic-augmented linear dot product} (Table~\ref{tab:miss}, last column).
  Adding a learned $-\frac12 v^\top Sv$ should fix P1 one-way. Adding $-\frac12 u^\top Ru$ as well
  should make two-way and joint exact, so P4 disappears. This answers open question~1.
\item \textbf{Weak-dependence Gaussian} (Prop.~\ref{prop:scalar}). On a domain with weak
  past-future dependence (scalar $\rho\approx0.3$, or a GP with a gap between past and future),
  trained on a lot of data, linear two-way/joint \Ltwo{} should trail one-way by roughly
  20--25\% of the information. The energy-score gap will be smaller but should have the same sign.
\item \textbf{Data-size sweep on weather} (\S\ref{sec:tension}). If the bias--variance reading
  is right, the gap between one-way and joint \Ltwo{} shrinks as training data grows.
\item \textbf{Normalised-cosine information cap} (Prop.~\ref{prop:twosided}). The learned
  forecast's $\KL(\nu_\vu\|\mu_v)$, which can be estimated from the bank weights, should never
  exceed $2\kappa$. Tuning $\kappa$ after training should control its calibration on its own
  (open question~2).
\end{enumerate}

%======================================================================
\section{Scope and caveats}

\begin{itemize}
\item Everything in Part~II is at the population level: the best member of each family under
  each loss, with infinite data and infinite negatives. Finite-sample statements are limited
  to Propositions~\ref{prop:snis} and~\ref{prop:collapse}, and to the hypothesis in
  \S\ref{sec:tension}.
\item The numbers in Table~\ref{tab:miss} are information (KL) losses. The experiments report
  the energy score, which is less sensitive to spread errors. The two rankings should agree
  in sign, but the sizes are not comparable.
\item Table~\ref{tab:miss} assumes full embedding dimension. The rank-constrained case (the
  paper's Thms~5.1, 5.3 and~5.6 with $r<\min(n_u,n_v)$) adds a further, separate source of
  misspecification.
\item Proposition~\ref{prop:relu} covers ReLU-type (piecewise-linear) networks. Other
  activations need their own growth bound.
\end{itemize}

%======================================================================
\appendix
\section{Proofs}

\paragraph{Proposition~\ref{prop:bounded}.} $Z(u)=\int e^{s}\,d\mu_v\le e^c\mu_v(V)=e^c$, and the
same bound holds for $Z$. For \Ltwo{}, $-\frac\kappa2|f-g|^2\le 0$. For normalised cosine,
Cauchy--Schwarz gives $|\langle\bar f,\bar g\rangle|\le1$. \qed

\paragraph{Proposition~\ref{prop:linear}.} $Z(u)=e^{b(u)}\E\exp\langle H^\top a(u),v\rangle$ by
definition. A multivariate $t$ is elliptical, so for $w\ne0$ the projection $\langle w,v\rangle$ is a
scaled univariate $t$ with $\nu$ degrees of freedom. Its density decays like $|x|^{-\nu-1}$,
so $\E e^{\langle w,v\rangle}=\infty$. \qed

\paragraph{Proposition~\ref{prop:relu}.} Write $g=W_L\sigma(\dots\sigma(W_1v+b_1)\dots)+b_L$ with
$\sigma=\mathrm{ReLU}$, and $g_\infty$ for the same map with all $b_\ell=0$. ReLU is
1-Lipschitz and positively homogeneous, so by induction over layers
$|g(v)-g_\infty(v)|\le\sum_\ell\|b_\ell\|\prod_{k>\ell}\|W_k\|=:B<\infty$, and $g_\infty(tv)=tg_\infty(v)$
for $t\ge0$. Let $\delta=\langle f(u),g_\infty(v_0)\rangle>0$. By continuity there is an open cone
$K\ni v_0$ on which $\langle f(u),g_\infty(v/|v|)\rangle\ge\delta/(2|v_0|)=:\delta'$. On $K$,
$s(u,v)\ge\kappa(\delta'|v|-|f(u)|B)$. An elliptical $t$ gives every open cone a tail
$\mu_v(K\cap\{|v|>t\})\asymp t^{-\nu}$, so $\int_K e^{\kappa\delta'|v|}\,d\mu_v=\infty$. \qed

\paragraph{Proposition~\ref{prop:snis}.} For bounded continuous $\varphi$,
$\frac1N\sum_j e^{s(u,v_j)}\varphi(v_j)\to\int e^{s}\varphi\,d\mu_v$ and
$\frac1N\sum_je^{s(u,v_j)}\to Z(u)$ a.s., by the strong law of large numbers, because both
integrands are $\mu_v$-integrable when $Z(u)<\infty$. Taking the ratio gives
$\int\varphi\,d\nu^N_\vu\to\int\varphi\,d\nu_\vu$. Applying the same law to the
countable convergence-determining class gives weak convergence a.s. For the ESS, write
$\rho=d\nu_\vu/d\mu_v=e^{s}/Z(u)$. Then
$N\sum_jw_j^2=\frac{\frac1N\sum_j\rho(v_j)^2}{(\frac1N\sum_j\rho(v_j))^2}\to\E_{\mu_v}[\rho^2]/1$.
If $Z(u)=\infty$, the law of large numbers for nonnegative variables gives
$\frac1N\sum_je^{s(u,v_j)}\to\infty$ a.s. \qed

\paragraph{Proposition~\ref{prop:collapse}.} Let $X_{(1)}\ge X_{(2)}\ge\cdots$ be the order
statistics and $b_N$ the $1-1/N$ quantile of $X$. Then $b_N=N^{1/\alpha}\ell(N)$ with $\ell$
slowly varying. By the classical extreme-value limit for regularly varying tails,
$(X_{(1)},X_{(2)})/b_N\Rightarrow(\Gamma_1^{-1/\alpha},\Gamma_2^{-1/\alpha})$, where
$\Gamma_1<\Gamma_2$ are the first two arrival times of a unit-rate Poisson process. These are a.s.\
distinct, so for every $\varepsilon>0$ there is $c>0$ with
$P\big(X_{(1)}-X_{(2)}\ge c\,b_N\big)\ge1-\varepsilon$ for large $N$. On that event,
\[
  1-\max_jw_j\;\le\;\sum_{j\ne(1)}e^{X_j-X_{(1)}}\;\le\;N\,e^{-c\,b_N}\;\to\;0,
\]
because $b_N$ grows polynomially and so faster than $\log N$. \qed

\paragraph{Proposition~\ref{prop:twosided}.} With $|s|\le\kappa$, we have $Z(u)\in[e^{-\kappa},e^\kappa]$,
so $d\nu_\vu/d\mu_v=e^{s}/Z(u)\in[e^{-2\kappa},e^{2\kappa}]$. Then
$\KL(\nu_\vu\|\mu_v)=\E_\nu\log(d\nu/d\mu_v)\le2\kappa$ and
$\E_{\mu_v}[(d\nu/d\mu_v)^2]=\E_\nu[d\nu/d\mu_v]\le e^{2\kappa}$, and the ESS limit follows from
Proposition~\ref{prop:snis}. \qed

\paragraph{Proposition~\ref{prop:target}.} (i) $J_\vu(s)=0$ iff
$\nu_\vu(\cdot\mid u)=\mu_\vu(\cdot\mid u)$ for a.e.\ $u$. Comparing densities with respect to
$\mu_v$, that means $e^{s(u,v)}/Z(u)=r(u,v)$, i.e.\ $s=\log r+\log Z(u)$. Conversely, any
$s=\log r+c(u)$ gives $Z(u)=e^{c(u)}\int r(u,\cdot)\,d\mu_v=e^{c(u)}$, and hence the true
conditional. (ii) is symmetric. (iii) With $\lambda_u,\lambda_v>0$, $J_2=0$ iff both (i) and
(ii) hold, so $c(u)=c'(v)$ for $\mu_u\otimes\mu_v$-a.e.\ $(u,v)$. A function of $u$ alone that
equals a function of $v$ alone almost everywhere under a product measure is a.e.\ constant.
For the joint loss, $\KL(\mu\|\nu)=0$ iff $e^{s}/Z=r$, i.e.\ $s=\log r+\log Z$. \qed

\paragraph{Proposition~\ref{prop:l2joint}.} The joint model has density proportional to
$e^{s}\,\mathcal N(0,D)$. With $s=-\frac12z^\top Qz$ and $z=(u,v)$, that is $\mathcal N(0,(D^{-1}+Q)^{-1})$.
Since $D^{-1}+Q\succeq D^{-1}$, inverting reverses the order and $(D^{-1}+Q)^{-1}\preceq D$. On
the data side, $D-C=\begin{psmallmatrix}0&-C_{uv}\\-C_{uv}^\top&0\end{psmallmatrix}$, whose
eigenvalues are $\pm\sigma_i(C_{uv})$. So $C\preceq D$ forces $C_{uv}=0$. Zero two-way or joint
loss requires the model joint to equal $\mu$ (Proposition~\ref{prop:target}(iii)), hence
$(D^{-1}+Q)^{-1}=C$, which is impossible when $C_{uv}\ne 0$. The Schur-complement inequality
$S\succeq A^\top R^+A$ is the generalised Schur criterion for
$\begin{psmallmatrix}R&-A\\-A^\top&S\end{psmallmatrix}\succeq0$. \qed

\paragraph{Proposition~\ref{prop:scalar}.} With $s=-\frac12(gu-hv)^2$ and unit variances,
$Q=\begin{psmallmatrix}g^2&-gh\\-gh&h^2\end{psmallmatrix}$ and, up to constants,
$2J_{\rm joint}=\mathrm{tr}((I+Q)C)-\log\det(I+Q)=g^2+h^2-2gh\rho-\log(1+g^2+h^2)$. For fixed
$t=g^2+h^2$ the cross term is minimised at $g=h$, which leaves $t(1-\rho)-\log(1+t)$, minimised at
$t=\rho/(1-\rho)$. The model forecast has precision $1+h^2=1+t/2=(2-\rho)/(2(1-\rho))$ and
mean $gh/(1+h^2)=\rho/(2-\rho)$. Expanding $J_\vu$ for small $\rho$ gives
$J_\vu=\frac{3}{16}\rho^2+O(\rho^3)$, while $I(u;v)=-\frac12\log(1-\rho^2)=\frac12\rho^2+O(\rho^4)$,
so the ratio tends to $3/8$. Table~\ref{tab:miss} confirms the closed form numerically: at
$\rho=0.5$ it gives $0.1699$. \qed

\begin{thebibliography}{9}
\bibitem{bst} R.~Baptista, A.~Stuart and S.~Tran.
  \emph{A Mathematical Perspective on Contrastive Learning}. arXiv:2505.24134, 2025.
\bibitem{cd} S.~Chatterjee and P.~Diaconis.
  The sample size required in importance sampling.
  \emph{Annals of Applied Probability}, 28(2):1099--1135, 2018.
\end{thebibliography}

\end{document}
